\chapter*{Preface}
Tome~5 is where the first four tomes meet. The integers of Tome~1, the characters of Tome~2, the probability of Tome~3
and the analysis of Tome~4 are used together to answer the oldest questions about prime numbers: how many there are
up to $x$, how they are shared among the residue classes modulo $q$, how they are spaced, and how the prime factors
of a typical integer are distributed. The summit is the prime number theorem, $\pi(x)\sim\frac x{\log x}$, proved in
Chapter~5 with Newman's short analytic argument, and then, with the zeros of the zeta function, with an error term
and Riemann's explicit formula. Around it stand Mertens's theorems, Dirichlet's theorem on primes in progressions,
sieves, the Erd\H{o}s--Kac theorem (the number of prime factors of an integer is asymptotically Gaussian), random
models of the primes, exponential sums, and a first taste of modular forms.

It assumes Tomes~0 to~4; references such as ``Tome~4, Theorem~10.3'' point there. Most of the analysis comes from
Tome~4: Abel summation (Chapter~1), the Euler--Maclaurin formula (Chapter~11), the Gamma function (Chapter~5),
Fourier analysis (Chapters~6--8), and complex analysis and the zeta function (Chapters~9 and~10). Characters of
finite abelian groups come from Tome~2, Chapter~11, and the central limit theorem and the method of moments from
Tome~3. In this tome $\log$ is the natural logarithm, as is usual in number theory, and $p$ always denotes a prime.

Exercises are graded \textbf{$\star$} warm-up, \textbf{$\star\star$} for a good afternoon, \textbf{$\star\star\star$}
gems, and every exercise has a hint and a complete solution at the end of its chapter. Each chapter has a
companion Python notebook. In this tome the computer is a laboratory: we count primes up to $10^9$, race residue
classes, rebuild the prime staircase from the zeros of $\zeta$, and test the random models against the data, and
the notebooks follow the ideas into computer science, from fast sieves to the testing of random number generators.

The rules of the earlier tomes hold here too. Every statement is proved, or, if quoted only for culture, labelled
as a remark that says where it is proved. Every solution starts with \emph{How to find it} before the
\emph{Write-up}. ``Pause and think'' boxes ask you to guess before you read on. The only fact used without proof is
the one already admitted in Tome~3, the Fubini--Tonelli theorem for general integrands; each use is labelled
\emph{Admitted (proved in Tome~6)}.

The chapters follow the plan of the series. Chapters~1 and~2 are elementary estimates: Mertens's theorems and the
average size of arithmetic functions. Chapters~3 and~4 prove Dirichlet's theorem. Chapters~5--7 prove the prime
number theorem, study the zeros of $\zeta$ and the explicit formula. Chapters~8--12 are sieves, probabilistic number
theory, random models, exponential sums, and modular forms.
